%!TEX Root = ../main.tex
\chapter{Integral Extensions}
\label{chap:integral}
In this chapter, we study a integral dependence
(being a root of monic polynomial),
which is a fundamental notion that unifies different
meanings in number theory and algebraic geometry.
\begin{itemize}
    \item In number theory, this is used to define
    algebraic integer, an analogue of $\ZZ$
    when the fraction field $\QQ$ is replaced 
    with a finite extension.
    \item In algebraic geometry, this is related to 
    (ramified) covering maps and resolving singularities
    of codimension~$1$.
\end{itemize}
We will explore various applications, especially 
in algebraic geometry.
\begin{itemize}
\item First, we will prove Hilbert Nullstellensatz
(\Cref{thm:max_ideal_poly_ring}), 
which classifies all maximal ideals of $\CC[x_1,\dots,x_n]$
as promised in the first chapter.
\item Second, in \Cref{sec:dim_intro},
we will explore how to define \textbf{dimension} of varieties
and prove some basic properties using tools we have so far
(although we will give a more robust treatment
of dimension again in \Cref{chap:dimension}).
\end{itemize}
\section{Integral Elements}
\begin{definition}[Integral Elements]
    Let $A\subset B$ be ring.
    We say that an element $b\in B$ is
    \vocab[integral element]{integral} over $A$ if
    $b$ is a root of monic polynomial
    $$b^n + a_{n-1}b^{n-1}  +\dots + a_0=0$$
    where $a_i\in A$.
\end{definition}
If $A$ and $B$ are fields, then the above definition
is equivalent to that $b$ is algebraic over $A$.

In the number-theoretic world, if $K$
is a number field (finite extension of $\mathbb Q$),
then the integral elements in $K$ form
the \vocab{ring of integers} $\mathcal O_K$,
which is an analogue of integers.
For example, if $K=\mathbb Q(\sqrt{-3})$,
we have $\calO_K = \mathbb Z\left[\frac{1+\sqrt{-3}}2\right]$.
\notation{$\calO_K$}{ring of integers}

We will now give a clean characterization of integrality.
To do so, we give the following intermediate definition.
\begin{definition}
    Given $b\in B$, define the subring generated by $b$
    by $A[b]$ being the smallest subring of $B$ containing $b$.
    \notation{$A[b]$}{smallest subring containing $B$}
\end{definition}
Now, we state the theorem.
\begin{theorem}[Characterization of Integral Elements]
    \label{thm:integral_elements}
    The following are equivalent.
    \begin{enumerate}[label=(\alph*)]
	\item $b$ is integral over $A$.
	\item $A[b]$ is finitely-generated $A$-module.
	\item There exists a ring $C$, which is 
	    a finitely-generated $A$-module, such that
	    $A\subset C\subset B$ and $b\in C$.
    \end{enumerate}
\end{theorem}
As an example, $\frac 13\in\mathbb Q$,
but $\mathbb Z\left[\frac 13\right]$ is not
finitely-generated as $\mathbb Z$-module since $\frac 13$
is not integral.
\begin{proof}
    \textbf{(a) $\Rightarrow$ (b).}
    Let $b$ is integral over $a$,
    meaning that for some $a_0,\dots,a_{n-1}\in A$, we have
    $b^n+a_{n-1}b^{n-1}+\dots+b_0=0$.
    This means that $b^n\in\Span(b^{n-1},\dots,b,1)$.
    Now, notice that
    $$b^k\in\Span(b^{n-1},\dots,b,1)
    \implies b^{k+1}\in\Span(b^n,\dots,b^2,b)
    \subset \Span(b^{n-1},\dots,b,1),$$
    so by induction, $b^k\in\Span(b^{n-1},\dots,b,1)$,
    so $A[b]$ is generated by $b^{n-1},\dots,b,1$, done.

    \textbf{(b) $\Rightarrow$ (c).} Take $C=A[b]$.

    \textbf{(c) $\Rightarrow$ (a).}
    The multiplication by $b$ (in $C$)
    can be represented as a matrix $T\in\Mat_m(A)$.
    By Cayley-Hamilton theorem,
    $p_T=0$, where $p_T$ is a characteristic polynomial of $T$.
    Thus, we have $p_T(1)=0$,
    so writing coefficients out gives
    $b^n + a_{n-1}b^{n-1} + \dots + a_1b+a_0=0$,
    as desired.
\end{proof}
\begin{corollary}[Integral elements form a subring]
    Let $A\subset B$ be rings.
    Then, 
    $$A' = \{b\in B\text{ integral over }A\}$$
    is a subring of $B$
\end{corollary}
\begin{proof}
    It suffices to show that $A'$ is closed under addition
    and multiplication.
    Let $r,s\in A'$,
    so by condition (b) in \Cref{thm:integral_elements},
    we get that $A[r]$ and $A[s]$ are finitely-generated.
    We have to show that $r+s$ and $rs$ are integral over $A$.
    
    To do this, we use condition (c) in \Cref{thm:integral_elements}.
    Since $A[r]$ and $A[s]$ are finitely-generated,
    we get that $A[r]\otimes_A A[s]$
    is finitely-generated.
    However, there is a surjective map
    $A[r]\otimes_A A[s]\to A[r,s]$,
    so $A[r, s]$ is finitely-generated.
    
    Since $r+s$ and $rs$ are both in $A[r,s]$,
    condition (c) is satisfied, so both $r+s$
    and $rs$ are integral over $A$, and we are done.
\end{proof}
\begin{remark}
    It is tempting to argue the following:
    $A[r+s]\subseteq A[r,s]$, so $A[r+s]$ is finitely-generated,
    so condition (b) in \Cref{thm:integral_elements} apply.

    This is incorrect because submodule of finitely-generated
    modules is not necessarily finitely-generated
    (but this would work if we assume that $A$ is Noetherian though).
    This is why we have to develop condition (c)
    in \Cref{thm:integral_elements}.
\end{remark}
%
%
%
%
%
\section{Integral Extension and Closure}
\subsection{Definition}
\begin{definition}[Integral extension]
    Let $A$ be a ring. We say that ring $B\supseteq A$
    is an \vocab{integral extension}
    if every element $b\in B$
    is integral over $A$.

    Similarly (but less-frequently used),
    the map $\phi : A\to B$ is \vocab[integral morphism]{integral}
    if $\phi(A)\subset B$ is an integral extension.
\end{definition}
\begin{example}
    \begin{itemize}
    \item $\ZZ[\sqrt{-3}]$ is an integral extension of $\ZZ$
    because $\sqrt{-3}$ is integral in $\ZZ$
    (it is a root of $x^2+3$).
    \item $\CC[t]$ is an integral extension of $\CC[t^2]$
    because $t$ is a root of polynomial $x^2-t^2$.
    \end{itemize}
\end{example}
\begin{definition}[Integral closure and integrally closed]
    Let $A\subset B$. Then, the ring $\overline A$ of all
    integral elements in $B$ is called the
    \vocab{integral closure} of $A$ inside $B$.

    If $\overline A=A$, then $A$ is \vocab{integrally closed} inside $B$.

    In the definition above,
    when $A$ is a domain and $B$ is omitted,
    we mean $B=\Frac A$.
    In other words, a domain $A$ is \vocab{integrally closed}
    (or \vocab{normal})
    if and only if $A$ is integrally closed inside
    $\Frac A$. 
\end{definition}
Let's do some examples.
\begin{example}[$\ZZ$ is integrally closed]
    $\ZZ$ is integrally closed (over $\QQ$). To see why,
    pick $\tfrac rs\in\QQ$ integral over $\ZZ$ where , i.e.,
    $$\left(\tfrac rs\right)^n + a_{n-1}\left(\tfrac rs\right)^{n-1} + \dots + a_1\left(\tfrac rs\right)
    + a_0 = 0$$
    for $a_0,a_1,\dots,a_{n-1}\in\ZZ$. 
    By dividing out common factor, assume $\gcd(r,s)=1$.
    Multiplying by $s^n$ both sides gives
    $$r^n =  - a_{n-1}r^{n-1}s + \dots + a_1rs^{n-1}
    + a_0s^n$$
    Then, $s$ divides the right hand side, so $s\mid r^n$,
    so $s=\pm 1$, implying that $\tfrac rs\in\ZZ$.
\end{example}
The argument in the above example generalizes directly to the following.
\begin{proposition}
    \label{prop:ufd_integrally_closed}
    UFD is integrally closed.
\end{proposition}
Thus, for example, $\ZZ[i]$ and 
$\CC[x_1,\dots,x_n]$ is integrally closed.

Let's see how integral closure can break
in both number-theoretic and geometric world.
First, let's do an NT example.
\begin{example}[Non-example in NT]
    $\mathbb Z[\sqrt{-3}]$ is not integrally closed
    since $\frac{1+\sqrt{-3}}2$ is integral.
    Thus, by \Cref{prop:ufd_integrally_closed}
    $\ZZ[\sqrt{-3}]$ is not a UFD.

    More generally, for any number field $K$, any UFD subring
    of $K$ must consists of only integral elements,
    and so it must be contained in the ring of integral element
    $\calO_K$. We will show later by \Cref{cor:integral_closure_one_step} that $\calO_K$ is integrally closed.
    This is one reason why $\calO_K$ is of a special interest
    in algebraic number theory.
\end{example}
Now, let's see how integral closure breaks in geometric world.
\begin{example}[Non-example in AG]
    \label{ex:integral_closure}
    We let 
    $$A = \frac{\mathbb C[x,y]}{(y^2-x^3)},$$
    the coordinate ring of a cuspidal curve.
    Let $u = \frac yx \in \Frac A$.
    We have $u^2 = \frac{y^2}{x^2} = \frac{x^3}{x^2} = x$,
    so $u$ is in the integral closure.

    Now, we consider the map
    \begin{align*}
	A &\to \mathbb C[u] \\ 
	x &\mapsto u^2 \\
	y &\mapsto u^3
    \end{align*}
    This map is well-defined and is injective, so $A=\mathbb C[u^2,u^3]$.
    We then see that $\mathbb C[u]$ is an integral closure.
\begin{center}
\begin{asy}
size(4cm,0);
import graph;
draw((-0.5,-1)--(-0.5,1));
real f(real x){
    return x^(3/2);
}
path p = graph(f,0,1);
draw(shift(0.5,0)*p);
draw(shift(0.5,0)*reflect((0,0),(1,0))*p);
label("$\implies$", (0,0));
label("$\operatorname{Spec}\mathbb C[u]$", (-0.5,-1), S);
label("$\operatorname{Spec}\frac{\mathbb C[x,y]}{(y^2-x^3)}$", (1,-1), S);
\end{asy}
\end{center}
Geometrically, we have $A\hookrightarrow\mathbb C[u]$,
inducing a map $\Spec\mathbb C[u]\to\Spec A$.
We have $\Spec\mathbb C[u]$ is a line,
but $\Spec A$ is a non-smooth curve with a cusp at $(0,0)$.
Thus, essentially,
integral closure is a process of making a ring smoother.
\end{example}


This process holds for not just $y^2=x^3$, but
for any ``one-dimensional'' complex varieties
(i.e., curves).
The coordinate ring of a variety is
integrally closed if and only if 
the curve is a smooth manifold.
If the curve is not smooth, then taking integral closure
resolves singularity by \emph{straightening} the curve.
What's surprising is that this process is canonical,
which is a special property of one-dimensional rings.

On the other hand, integral closure does not resolve 
singularities for higher-dimensional varieties.
For example, in \Cref{prob:integrally_closed},
you will be asked to show that $\CC[x,y,z]/(xy-z^2)$
is integrally closed, despite that the surface 
$xy-z^2=0$ has a singularity at $(0,0,0)$.
In higher dimension, it is much more difficult to
resolve singularities, and often there are choices involved.
In \Cref{chap:dvr}, we will see how being integrally closed
is connected to smoothness.
\begin{center}
\begin{asy}
settings.tex="pdflatex";
settings.prc=false;
settings.render=0;
size(4cm,0);
import three;
import graph3;

currentprojection = orthographic(0,10,3);
path3 c1 = scale(2,1,1)*circle((0,0,-2),1,Z);
path3 c2 = scale(2,1,1)*circle((0,0,2),1,Z);

draw(extrude(c1, (0,0,0)--cycle), material(gray,emissivepen=gray));
draw(extrude(c2, (0,0,0)--cycle), material(gray,emissivepen=gray));
dot((0,0,0), red);
\end{asy}
\end{center}
\begin{center}
Surface $xy-z^2=0$.
\end{center}
%
%
%
%
%
\subsection{Basic Properties}
We now start to present properties of integral extension.
Before that, we make a remark about finitely-generated algebras. 
\begin{definition}[Finitely-generated vs. finite]
    Given a ring homomorphism $\phi : A\to B$.
    We say that 
    \begin{itemize}
    \item $B$ is a \vocab[finitely-generated algebra]
    {finitely-generated}
    (or \vocab[finite-type algebra]{finite-type}) $A$-algebra
    if and only if there is a surjection from
    $A[t_1,\dots,t_n]\to B$.
    \item $B$ is a \vocab[finite algebra]{finite} $A$-algebra
    if and only if $B$ is finitely-generated as $A$-module.
    \end{itemize}
\end{definition}
\begin{warn}
Don't confuse between finitely-generated and finite.
For example, polynomial ring $A[x]$ 
is a finitely-generated $A$-algebra but not a finite $A$-algebra.
\end{warn}
\begin{corollary}
    \label{cor:integral_finite_generation}
    Let $B$ be an $A$-algebra.
    Suppose that $B$ is finitely-generated
    as $A$-algebra and is integral over $A$,
    then $B$ is finitely-generated as an $A$-module.
\end{corollary}
\begin{proof}
    Follows directly from \Cref{thm:integral_elements}.
\end{proof}
\paragraph{Tower of Integral Elements}
\begin{proposition}[Integrality is transitive]
    \label{prop:integral_transitive}
    Let $A\subset B\subset C$. Suppose $B$ is integral over $A$
    and $C$ is integral over $B$, then
    $C$ is integral over $A$.
\end{proposition}
\begin{proof}
    Let $\gamma\in C$. Then, for some $b_0,\dots,b_{n-1}\in B$,
    we have $\gamma^n + b_{n-1}\gamma^{n-1} + \dots + b_1\gamma+b_0=0$.
    Let $B'=A[b_0,\dots,b_{n-1}]$,
    which is a finitely-generated $A$-algebra,
    so by \Cref{cor:integral_finite_generation},
    it is a finitely-generated $A$-module.
    Since $B'[\gamma]$ is a finitely-generated $B'$-module, it is a finitely-generated $A$-module,
    so by \Cref{thm:integral_elements} ((c) implies (a)), $\gamma$ is integral over $A$.
\end{proof}
\begin{corollary}[Integral closure is integrally-closed]
    \label{cor:integral_closure_one_step}
    Let $A\subset B$ be rings.
    If $\overline A$ is an integral closure of $A$ inside $B$,
    then $\overline A$ is integrally closed already.
\end{corollary}
\begin{proof}
    If $\gamma\in B$ is integral over $\overline A$,
    then by the proposition above,
    $\gamma$ is integral over $A$, so $\gamma\in\overline A$.
\end{proof}
\paragraph{Integral Closure vs. Quotient and Localization}
\begin{proposition}[Integral extension is preserved
    under quotient and localization]
    \label{prop:integral_quotient_localize}
    Let $A\subseteq B$ be rings such that $B$ 
    is integral over $A$.
    \begin{alphenum}
    \item If $J$ is an ideal of $B$ and 
    $I = J\cap A$, then $B/J$ is integral over $A/I$.
    \item If $S\subset A$ be multiplicatively
    closed, then $S^{-1}B$ is integral over $S^{-1}A$.
    \end{alphenum}
\end{proposition}
\begin{proof}
    \begin{alphenum}
    \item Any $b\in B$ satisfies the relation 
    $b^n + a_{n-1}b^{n-1} + \dots + a_1b +a_0=0$
    for some $a_0,a_1,\dots,a_{n-1}\in A$.
    Now reduce this equation modulo $J$.
    \item Note that $\frac bs = \frac 1s\cdot b$, and we have
    $\frac 1s$ and $\frac b1$ are both integral over $S^{-1}A$.
    \qedhere
    \end{alphenum}
\end{proof}
\begin{proposition}[Localization preserves integral closure]
\label{prop:integral_closure_local}
Let $A\subseteq B$ be rings,
and let $C$ be the integral closure of $A$ in $B$
(so $A\subseteq C\subseteq B$).
Then, $S^{-1}C$ is the integral closure of $S^{-1}A$
in $S^{-1}B$.
\end{proposition}
\begin{proof}
    By \Cref{prop:integral_quotient_localize} (b),
    we see that every element in $S^{-1}C$ is integral
    over $S^{-1}B$. 
    Conversely, if $\tfrac bs\in S^{-1}B$ is integral in
    $S^{-1}A$, then we have the equation in form
    $$\left(\frac bs\right)^n 
    + \frac{a_{n-1}}{s_{n-1}}\left(\frac bs\right)^n
    + \dots + \frac{a_1}{s_1}\left(\frac bs\right)
    + \frac{a_0}{s_0} = 0,$$
    where $a_0,\dots,a_{n-1}\in A$ and $s_0,\dots,s_{n-1}\in S$.
    Let $t = s_0\cdots s_{n-1}$. 
    Then, multiplying this equation by $(ts)^n$ gives 
    $$(bt)^n + c_{n-1}(bt)^{n-1} + \dots + c_1(bt) + c_0=0$$
    for some $c_0,\dots,c_{n-1}\in A$.
    This shows that $bt\in C$, so $b/s = bt/st \in S^{-1}C$.
\end{proof}
Finally, we note that being integrally closed is a stalk-local property.
\begin{proposition}[Being integrally closed is stalk-local]
    \label{prop:integrally_closed_local}
    If $A$ is a \textbf{domain}, then the following are equivalent.
    \begin{alphenum}
    \item $A$ is integrally closed.
    \item $A_\pp$ is integrally closed for all prime ideal $\pp$.
    \item $A_\mm$ is integrally closed for all maximal ideal $\mm$.
    \end{alphenum}
\end{proposition}
\begin{proof}
    \Cref{prop:integral_closure_local} gives 
    that (a) implies (b), and clearly, (b) implies (c).
    Thus, it suffices to prove that (c) implies (a).

    Assume (c). Let $B$ be the integral closure of $A$,
    and let $\phi : A\to B$ be the inclusion map.
    Then, by \Cref{prop:integral_closure_local},
    $B_\mm$ is the integral closure of $A_\mm$,
    so the map $\phi_\mm : A_\mm \to B_\mm$
    is surjective for all maximal ideals $\mm$.
    Since being surjective is local (\Cref{cor:injective_surjective_local}),
    we get that $\phi$ is surjective, so $A=B$.
\end{proof}
\paragraph{The Minimal Polynomial Test}
Verifying that $\alpha$ is not integral over $A$
would be difficult since we need to examine all
possible monic polynomials vanishing at $\alpha$.
Fortunately, with favorable condition,
we only need to examine the minimal polynomial.
\begin{lemma}[Minimal Polynomial Test]
    \label{lem:min_poly_test}
    Let $A$ be an integrally-closed domain.
    Let $K=\Frac(A)$, and let $L$ be a field extension of $K$.
    Let $\alpha\in L$, and let $f(x)\in K[x]$ be the monic minimal polynomial of $\alpha$ over $K$.

    Then, $\alpha$ is integral over $A$ if and only if
    $f(x)\in A[x]$.
\end{lemma}
\begin{proof}
    If $f(x)\in A[x]$, then $\alpha$ is clearly integral over $A$ by definition.

    In another direction, let $g(x)\in A[x]$ be a monic polynomial such that $g(\alpha)=0$.
    Let $\alpha=\alpha_1, \alpha_2,\dots,\alpha_n$
    be roots of $f$ in $\overline K$.
    Thus, since $f$ divides $g$ in $K[x]$,
    $\alpha_i$ must also be a root of $g$, 
    and hence integral over $A$.
    Therefore, each coefficient of polynomial
    $$f(x) = (x-\alpha_1)\cdots (x-\alpha_n)$$
    is a sum of products of $\alpha_i$, and hence integral
    over $A$. However, $A$ is integrally closed, so
    all coefficients are in $A$, or $f\in A[x]$.
\end{proof}
\begin{example}
Using the above \Cref{lem:min_poly_test}
with $A=\ZZ$ and $K=\QQ$, we can see that
$\tfrac{1+\sqrt{3}}2$ is not integral over $\ZZ[\sqrt 3]$
because the minimal polynomial $x^2-x-\tfrac 12$
does not have integer coefficients.
\end{example}

\section{Cohen-Seidenberg Theory}
Cohen-Seidenberg theory is a series of theorems
that control prime ideals in an integral extensions,
making them useful in algebraic geometry.
In this section, we describe and prove those.
\subsection{Going-up Theorem}
\paragraph{Lying Over}
First, let us define a new terminology.
\begin{definition}[Lies over]
    Consider an integral extension $A\subset B$.
    A prime ideal $\mathfrak q\subset B$ \vocab{lies over}
    a prime ideal $\mathfrak p\subset A$
    if and only if $\mathfrak p = \mathfrak q\cap A$.

    Another interpretation: the inclusion map 
    $\phi : A\hookrightarrow B$ induces the map on spectrum
    $\phi^\# : \Spec B\to \Spec A$ taking $\qq$ to $\qq\cap A$.
    Thus, $\qq$ lies over $\pp$ if $\phi^\#\qq  = \pp$.
\end{definition}
\begin{example}[Gaussian integers]
    If $A=\ZZ$ and $B=\ZZ[i]$, then
    \begin{itemize}
    \item $(3)\subseteq \ZZ[i]$ is the only prime
    lying over $(3)\subseteq\ZZ$.
    \item $(2+i)$ and $(2-i)$ are two primes
    lying over $(5)\subseteq\ZZ$.
    \end{itemize}
    The map $\Spec \ZZ[i] \to \Spec\ZZ$ looks like 
    a \emph{double cover}, as in the following diagram
    \begin{center}
    \begin{asy}
    size(10cm);
    for(int i=1; i<=6; ++i){
        draw((i,0.5)--(i,3), dotted);
    }

    draw((0.5,0.5) -- (6.7,0.5), red);
    dot("$(2)$", (1,0.5), S, red);
    dot("$(3)$", (2,0.5), S, red);
    dot("$(5)$", (3,0.5), S, red);
    dot("$(7)$", (4,0.5), S, red);
    dot("$(11)$", (5,0.5), S, red);
    dot("$(13)$", (6,0.5), S, red);
    label("$\operatorname{Spec}\mathbb Z$", (6.7,0.5), E);

    dot("$(1+i)$", (1,2), W, blue);
    dot("$(3)$", (2,2), S, blue);
    dot("$(2+i)$", (3,2.5), N, blue);
    dot("$(2-i)$", (3,1.5), S, blue);
    dot("$(7)$", (4,2), S, blue);
    dot("$(11)$", (5,2), S, blue);
    dot("$(3+2i)$", (6,2.5), N, blue);
    dot("$(3-2i)$", (6,1.5), S, blue);

    draw((6.7,2.5){W} .. (6,2.5) .. (5,2.03){W} .. (4.5,2.3){W} .. (4,2.03){W}
    .. (3,2.5) .. (2,2.03){W} .. (1.5,2.3){W} .. (1,2)
    .. (1.5,1.7){E} .. (2,1.97){E} .. (3,1.5)
    .. (4,1.97){E} .. (4.5,1.7){E} .. (5,1.97){E} .. (6,1.5) .. (6.7,1.5){E}, blue);

    label("$\operatorname{Spec}\mathbb Z[i]$", (6.7,1.5), E);
    \end{asy}
    \end{center}
\end{example}
\begin{example}
Suppose that $A=\CC[t^2]$ and $B=\CC[t]$.
Then, the map $\Spec \CC[t]\to\Spec \CC[t^2]$
takes prime ideal $(t-a)$ to prime ideal $(t^2-a^2)$. 
This map correspond to the morphism 
on affine line $\CC\to \CC$ that sends $a$ to $a^2$.

Visually, the map $\Spec \CC[t] \to \Spec \CC[t^2]$
looks like a double cover map, ramified at $0$,
as in the following diagram.
\begin{center}
\begin{asy}
size(10cm);
defaultpen(fontsize(10pt));
import graph;

real f(real x){return 2*x^2;}
draw(rotate(-90)*graph(f,-sqrt(4.5),sqrt(4.5)), blue);

for(int i=0; i<=4; ++i){
	if(i>0){
		dot("$(x-\sqrt{" + ((string)i) + "})$", (2*i,sqrt(i)), N, blue);
		dot("$(x+\sqrt{" + ((string)i) + "})$", (2*i,-sqrt(i)), S, blue);
		dot("$(x^2-" + ((string)i) + ")$", (2*i,-3), S, red);
	}
	else{
		dot("$(x)$", (0,0), W, blue);
		dot("$(x^2)$", (0,-3), S, red);
	}
	draw((2*i,-3) -- (2*i,2.5), dotted);
}
draw((-0.5,-3)--(9,-3), red);
label("$\operatorname{Spec}\mathbb C[x^2]$", (9,-3), E);
label("$\operatorname{Spec}\mathbb C[x]$", (9,-sqrt(4.5)),E);
\end{asy}
\end{center}
\end{example}

Informally, if $B$ is an integral extension of $A$,
then the map $\Spec B\to \Spec A$ looks like a ramified 
covering map. 
The algebraic results that explain this 
are the following theorems.
\paragraph{Key Results}
We will now state three imporant theorems.
\begin{theorem}[Lying Over]
    \label{thm:lying_over}
    \index{Lying Over Theorem}
    Let $A\subset B$ be an integral extension.
    For any prime $\mathfrak p\subset A$,
    there exists a prime ideal $\mathfrak q\subset B$
    lying over $\mathfrak p$.
    \begin{center}
    \begin{tikzcd}
        {\color{gray}\qq}  \arrow[d, no head, dashed] \arrow[r,"\subset",phantom]
        & B \arrow[d, no head] \\
        \pp   \arrow[r,"\subset",phantom]
        & A 
    \end{tikzcd}
    \end{center}
    (In other words, the map $\Spec B\to\Spec A$ is surjective.)
\end{theorem}
\begin{theorem}[Incomparability]
    \label{thm:lying_over_equal}
     Let $A\subseteq B$ be an integral extension,
    $\pp\subset A$ be a prime ideal.
    If $\qq\subseteq\qq'$ be two prime ideals
    lying over $\pp$, then $\qq=\qq'$.
\end{theorem}
\begin{theorem}[Going-up]
    \label{thm:going_up}
    \index{Going-up Theorem}
    Let $A\subseteq B$ such that $B$ is an
    integral extension of $A$.
    Given two prime ideals
    $\mathfrak p\subset\mathfrak p'\subset A$,
    and another prime ideal $\mathfrak q\subset B$
    such that $\mathfrak q\cap A = \mathfrak p$.
    Then, there is another prime ideal $\mathfrak q'$
    such that 
    $\qq\subset\qq'\subset B$
    and $\mathfrak q'\cap A = \mathfrak p'$.
    \begin{center}
    \begin{tikzcd}
        \qq  \arrow[d, no head] \arrow[r,"\subset",phantom]
        & {\color{gray}\qq'}  \arrow[d, no head, dashed] \arrow[r, "\subset",phantom]
        & B \arrow[d, no head] \\
        \pp   \arrow[r,"\subset",phantom]
        & \pp'  \arrow[r, "\subset",phantom]
        & A 
    \end{tikzcd}
    \end{center}
\end{theorem}
\begin{exercise}[Morphism version of the above theorems]
    Let $\phi : A\to B$ be an integral morphism.
    Using the above theorems, prove the following.
    \begin{alphenum}
    \item For any prime ideal $\pp\subset A$,
    prove that there exists a prime ideal $\qq\subset B$
    such that $\phi^{-1}\qq = \pp$.
    \item Given two primes ideal $\pp\subset \pp'\subset A$
    and another prime ideal $\qq\subset B$ such that 
    $\phi^{-1}\qq = \pp$,
    prove that there exists a prime ideal $\qq'$ such that
    $\qq\subset\qq'\subset B$ and $\phi^{-1}\qq' = \pp'$.
    \end{alphenum}
\end{exercise}
\begin{exercise}
    Deduce from going-up theorem that 
    if $\phi : A\to B$ is integral map,
    then $\phi^\# : \Spec B\to\Spec A$ is a \vocab{closed map}
    (i.e., image of closed sets is closed).
\end{exercise}
\paragraph{Proof}
The rest of this subsection is devoted to proving these theorems.
We start with a very special case
when $B$ is a field.
\begin{lemma}[The field case]
    \label{lem:field_case}
    Let $A\subseteq B$ such that $A$ and $B$ are domains and $B$ is an integral extension of $A$.
    Then, $A$ is a field if and only if $B$ is a field.
\end{lemma}
\begin{proof}
    ($\Rightarrow$) 
    Let $b\in B$. There exists $c_0,\dots,c_{n-1}\in A$
    such that 
    $$b^n + c_{n-1}b^{n-1} + \dots + c_1b + c_0 = 0.$$
    We may assume $c_0\neq 0$ because otherwise,
    one can divide out by $b$ (this requires the fact that 
    $B$ is a domain).
    Since $A$ is a field, there exists $d\in A$ such that 
    $dc_0=1$. Then, we have
    $$b\cdot -d(b^{n-1}+c_{n-1}b^{n-2} +\dots + c_2b+c_1)  = 1,$$
    giving the inverse of $b$.

    ($\Leftarrow$)
    Let $a\in A$. Since $B$ is a field, 
    we have $a^{-1}\in B$.
    Thus, for some $c_0,\dots,c_{n-1}\in A$, we have
    \begin{align*}
	(a^{-1})^n + c_{n-1}(a^{-1})^{n-1} + \dots + 
	c_1a^{-1} + c_0 &= 0 
    \intertext{Multiplying by $a^n$ both sides gives}
        1 + c_{n-1}a + \dots + c_1a^{n-1} + c_0a^n &= 0 \\
	1 - a\cdot(\text{fluff}) &= 0,
    \end{align*}
    and this fluff will be the inverse of $a$.
\end{proof}
Believe it or not, by solving a very specific case, 
we actually have half of the work done.
The strategy is to use quotient and localization
to reduce the theorem down to the above lemma.

\begin{corollary}[The maximal ideal case]
    \label{cor:max_ideal_case}
    Let $A\subseteq B$ be integral,
    $\pp\subset A$ be a prime ideal,
    and $\qq$ lies over $\pp$.
    Then, $\qq$ is maximal if and only if $\pp$ is maximal.
\end{corollary}
\begin{proof}
    Since integral closure is preserved under quotient (\Cref{prop:integral_quotient_localize} (a)),
    $B/\qq$ is an integral extension of $A/\pp$.
    Thus, by the field case (\Cref{lem:field_case}),
    $B/\qq$ is a field if and only if $A/\pp$ is a field.
    This implies the result.
\end{proof}

Now, we can prove the main theorems.
\begin{proof}[Proof of Lying Over \Cref{thm:lying_over}]
We use localization. As above, let 
$S=A\setminus \pp$ and $B_{\mathfrak p} = S^{-1}B$
and $B_\pp$ is integral over $A_\pp$ by 
\Cref{prop:integral_quotient_localize} (b).
Then, let $\mathfrak m$ be maximal ideal in $B_{\mathfrak p}$.
By the maximal ideal case \Cref{cor:max_ideal_case}, $\mathfrak m \cap A_{\mathfrak p}$ must be (the only) maximal ideal in $A_\pp$, which is be $\mathfrak p A_{\mathfrak p}$.
Thus, $\frakm \cap A_\frakp = \frakp A_\frakp$,
so taking intersection to $A$ both sides gives
$\frakm\cap A = \frakp$.
Hence, $\frakm\cap B$ is our desired ideal $\qq$.
\end{proof}

\begin{proof}[Proof of Incomparability 
\Cref{thm:lying_over_equal}.]
    We use localization.
    Let $S=A\setminus\mathfrak p$ and let
    $B_{\mathfrak p} = S^{-1}B$.
    Then, by \Cref{prop:integral_quotient_localize} (b),
    we get an integral extension 
    $A_{\mathfrak p}\subset B_{\mathfrak p}$.
    Now, if $\mathfrak q\subset\mathfrak q'$,
    then we extend to $B_{\mathfrak p}$
    to get $\mathfrak qB_{\mathfrak p}\subseteq
    \mathfrak q'B_{\mathfrak p}$.
    
    Now, the magic happens, since $\mathfrak pA_{\mathfrak p}$
    is a maximal ideal in $A_{\mathfrak p}$,
    we get by the maximal ideal case (\Cref{cor:max_ideal_case})
    that $\mathfrak qB_{\mathfrak p}$ 
    is a maximal ideal in $B_{\mathfrak p}$.
    This forces $\mathfrak qB_{\mathfrak p}
    = \mathfrak q'B_{\mathfrak p}$
    or $\mathfrak q = \mathfrak q'$, done.
\end{proof}

\begin{proof}[Proof of Going Up \Cref{thm:going_up}]
We use quotient again.
Since $\mathfrak q$ lies over $\mathfrak p$,
we have by \Cref{prop:integral_quotient_localize} (a) that $A/\mathfrak p \subset B/\mathfrak q$
is integral.
Thus, by lying over \Cref{thm:lying_over}, pick 
prime ideal $\overline{\mathfrak q'}\subset B/\mathfrak q$
lying over $\overline{\mathfrak p'}\subset A/\mathfrak p$.
By correspondence theorem,
we have $\mathfrak p\subset\mathfrak p'$
and $\mathfrak q\subset\mathfrak q'$.
Moreover, the lying over condition
$\overline{\mathfrak q'}\cap A = \overline{\mathfrak p'}$
directly translates to
$\mathfrak q'\cap A = \mathfrak p'$,
so we are done.
\end{proof}
% The last part of this theorem allows us to lift the chain
% $$\mathfrak p_1\subset\mathfrak p_2\subset\dots\subset
% \mathfrak p_r\subset A$$
% and a prime ideal $\qq_1$ to
% $$\mathfrak q_1\subset\mathfrak q_2\subset\dots
% \subset\mathfrak q_r\subset B.$$
% We will see later that prime ideal chains control the dimension,
% so this theorem controls the dimension of the ring.


% This theorem states that the map $\Spec B\to \Spec A$
% (where $A\to B$ integral) is a closed map.

\subsection{Going-down Theorem}
We will now state a similar result.
However, it will require additional assumption that
\begin{itemize}
    \item $B$ is a domain.
    \item $A$ is integrally closed.
\end{itemize}
\begin{remark}
Neither of these hypothesis can be dropped.
\Cref{prob:going_down} will ask you to translate the 
images to the counterexamples.
\end{remark}
\begin{theorem}[Going-down]
    \index{Going-down Theorem}
    \label{thm:going_down}
    Let $A\subset B$ be domains such that $A$ integrally closed
    and $B$ is integral over $A$.
    Given that $\mathfrak p\subset\mathfrak p'\subset A$
    be prime ideals and $\mathfrak q'$ lies over $\mathfrak p'$,
    then there exists a prime ideal $\mathfrak q\subset\mathfrak q'$
    lying over $\mathfrak p$.
    \begin{center}
    \begin{tikzcd}
        {\color{gray}\qq}  \arrow[d, no head, dashed] \arrow[r,"\subset",phantom]
        & \qq'  \arrow[d, no head] \arrow[r, "\subset",phantom]
        & B \arrow[d, no head] \\
        \pp   \arrow[r,"\subset",phantom]
        & \pp'  \arrow[r, "\subset",phantom]
        & A 
    \end{tikzcd}
    \end{center}
\end{theorem}
\begin{remark}
There is another variant of Going-down theorem
for flat modules, which will be stated and proved in 
\Cref{thm:going_down_flat}.
\end{remark}
To prove this theorem, we begin with a small lemma that allows us to find ideals.
\begin{lemma}[Lying Over Criterion]
    \label{lem:lying_over_criterion}
   Let $A\subset A'$ be a ring extension
   (not necessarily integral!).
   For any prime ideal $\pp\subset A$ the following are equivalent:
   \begin{alphenum}
       \item There exists a prime ideal $\pp'\subset A'$ such that $\pp'\cap A=\pp$.
       \item $\pp A' \cap A = \pp$.
   \end{alphenum}
\end{lemma}
\begin{proof}
    We first prove that (a) implies (b).
    To do this, we note that $\pp A'\subseteq\pp'$, so
    $$\pp \ \subseteq\ \pp A'\cap A \ \subseteq\ \pp'\cap A \ =\ \pp,$$
    so all equalities occur, implying that $\pp A'\cap A=\pp$.

    Next, we prove that (b) implies (a).
    The strategy is that we want prime ideal 
    that does not meet $A\setminus\pp$, so we invert that set.
    Localize $A'$ with multiplicative set 
    $S=A\setminus\pp$ to get $A_\pp$.
    Then, $\pp A' \cap S = (\pp A'\cap A)\setminus\pp =
    \emptyset$ (by assumption (b)). 
    Thus, $\pp A'_\pp$ is a proper ideal of $A'_\pp$,
    and hence contained in a maximal ideal $\mm \subseteq A'_\pp$.
    By \Cref{thm:prime_ideal_localize},
    $\mm$ must be in form $\pp' A'_{\pp}$ for some 
    prime ideal $\pp'\subset A'$ such that $\pp'\cap (A\setminus\pp)\neq\emptyset$.
    Since $\pp A'_\pp\subset \mm$, it follows that 
    $\pp\subset \pp'$.
    Combining previous two sentences gives $\pp'\cap A=\pp$.
\end{proof}
% In the world of Algebraic Geometry,
% $B$ is also a finitely-generated ring extension,
% then the map $\Spec B\to\Spec A$ is open map.
\begin{proof}[Proof of Going Down \Cref{thm:going_down}]
    First, we claim that it suffices to show
    $$\mathfrak p B_{\mathfrak q'}\cap A = \mathfrak p.$$
    Indeed, assume that this is true,
    then by \Cref{lem:lying_over_criterion}
    we there is a prime ideal of $B_{\qq'}$ lying above 
    $\pp A_\pp$, which must be of the form 
    $\qq B_{\qq'}$ for some prime ideal $\qq\subseteq \qq'$
    by \Cref{thm:prime_ideal_localize}.
    Since $\qq B_{\qq'}$ contracts to $\pp\subset A$,
    it follows that $\qq\subset B$ contracts to $\pp\subset A$.

    Now, we prove that $\mathfrak p B_{\mathfrak q'}\cap A
    = \mathfrak p$.
    One containment $\mathfrak p\subset \mathfrak pB_{\mathfrak q'}
    \cap A$ is automatic.
    Thus, suppose $a=\mathfrak pB_{\mathfrak q'}\cap A$
    but $a\notin\mathfrak p$. 
    Then, $a=\frac bs$, where $b\in\mathfrak p B$
    and $s\notin\mathfrak q'$.
    The proof proceeds in four steps.
    \begin{enumerate}
    \item \textbf{There exists polynomial $F(X)
    = X^n + c_{n-1}X^{n-1} + \dots + c_1X+c_0$
    such that $c_0,\dots,c_{n-1}\in\pp$ and $F(b)=0$.}
    The fact that $b\in\mathfrak pB$ means that
    $b$ can be written as a linear combination
    $b = \sum_{i=1}^m x_ip_i$, where $x_i\in B$ and $p_i\in \mathfrak p$.

    Set $A' = A\langle x_1,\dots,x_m\rangle\subset B$
    to be a finitely-generated submodule of $B$.
    We have $bA'\subset\mathfrak pA'$,
    so by Cayley-Hamilton theorem \Cref{cor:cayley_hamilton},
    there exists polynomial
    $$F(X) = X^n + c_{n-1}X^{n-1}+\dots+c_1X+c_0,\ \quad c_i\in\mathfrak p$$
    such that $F(b)=0$.
    \item \textbf{The minimal polynomial $G(X)=X^m + d_{m-1}X^{m-1} + \dots + d_1X+d_0$
    of $b$ satisfies $d_0,\dots,d_{m-1}\in \pp$.}
    Let $G$ be such minimal polynomial, so by the minimal
    polynomial test (\Cref{lem:min_poly_test}),
    we get that $G(X)\in A[x]$.
    Since $G(X)$ divides $F(X)$ in $K[X]$ and both polynoimals
    are monic, we get that the quotient must be in $A[x]$,
    so $F(X) = G(X)H(X)$ for some $H(X)\in A[X]$.

    Now, we reduce modulo $\pp$.
    We get $X^n = \overline G(X) \overline H(X)$
    in $A/\pp[x]$. Thus, $\overline G(X) = X^m$,
    so all other coefficients in $G$ vanish modulo $\pp$.
    \item \textbf{$d_i/a^{n-i} \in A$ for all $i=0,1,\dots,n-1$.}
    The minimal polynomial $F$ of $b$ can be translated
    to minimal polynomial of $s=\frac ba$ by
    \begin{equation}
    \label{eq:going_down}
    \left(\frac ba\right)^n
    + \frac{d_{n-1}}a \left(\frac ba\right)^{n-1}
    + \dots + \frac{d_0}{a^n} = 0.
    \end{equation}
    Thus, by the minimal polynomal test (\Cref{lem:min_poly_test}), we get that
    all coefficients of this polynomial must be in $A$, or
    $d_i/a^{n-i}\in A$ for all $i=0,1,\dots,n-1$.
    \item \textbf{Conclusion.}
    Let $e_i = d_i/a^{n-i}$.
    We have $d_i\in\pp$ from Step 2.
    However, $a\notin\pp$, so we have $e_i\in\pp$.
    
    Finally, we note that $\mathfrak p\subset\mathfrak q'$,
    so $e_i\in\mathfrak q'$ for all $i$.
    Thus, $e_is^i\in\mathfrak q'$ for all $i$.
    However, from \eqref{eq:going_down},
    we have $s^n+e_{n-1}s^{n-1}+\dots+e_1s+e_0=0$.
    All terms except $s^n$ are in $\qq'$, so $s^n\in\qq'$,
    and so $s\in\qq'$, a contradiction.
    \qedhere
\end{enumerate}
\end{proof}
%
%
%
%
%
\section{Noether Normalization and Nullstellensatz}
In this section, we discuss Noether Normalization Lemma,
an important source of integral extension that
can be obtained from any ring in form $k[x_1,\dots,x_n]/\pp$,
where $\pp$ is a prime ideal.
We will then use it to prove Hilbert's Nullstellensatz.
\subsection{Noether Normalization}
\begin{theorem}[Noether Normalization Lemma]
    \index{Noether Normalization Lemma}
    \label{thm:noether_normalization}
    Let $k$ be a field, and let $B$ be a
    finitely-generated $k$-algebra.
    In other words, $B=k[y_1,\dots,y_n]/I$
    for some ideal $I$.
    Then, there exists $x_1,\dots,x_m\in B$ such that
    \begin{itemize}
    \item The subring $A\subset B$ generated by
    $x_1,\dots,x_m$ is isomorphic to the polynomial ring
    $k[x_1,\dots,x_m]$
    (i.e., $x_1,\dots,x_m$ are algebraically independent).
    \item $B$ is an integral extension of $A$.
    \end{itemize}
\end{theorem}
In particular, the theorem states that ``any finitely-generated $k$-algebra is an integral extension of polynomial ring
$k[T_1,\dots,T_m]$.''
Geometrically, this means that any variety is 
a finite cover of the affine space $k^m$.

Before proving this theorem, we provide some example
of why this theorem is intuitive, but not completely obvious.
\begin{example}
Consider
$$B = \frac{k[x,y]}{(xy-1)}.$$
Let's randomly take $A=k[x]$, but this doesn't work
since one can check that $B$ is integral over $A$.
Geometrically, one can see that a portion of 
the hyperbola $xy=1$ got squeezed to near $x=0$,
making the map not look like covering map.

Thus, to fix this, we want to perturb $A$ a little bit.
In this case, we see that $A=k[x+\delta y]$
for any $\delta\neq 0\in k$ works. To see why, if $z=x+\delta y$,
then we have the relation $(z-\delta y)y = 0$,
which gives a monic polynomial that $y$ is a root.
Similarly, $(z-x)x = 0$, giving a monic polynomial that
$x$ is a root.
Geometrically, the projection map now looks like 
a $2$-to-$1$ covering map.
\begin{center}
\begin{adjustbox}{valign=m}
\begin{asy}
size(0,6cm);
import graph;
real f(real x){return 1/x;}
draw(graph(f,0.2,5), red+linewidth(1));
draw(graph(f,-0.2,-5), red+linewidth(1));
draw((-6,0)--(6,0), blue+linewidth(1));
for(int i=-4; i<=4; ++i){
    if(i!=0){
        draw((i,1/i)--(i,0), gray);
        draw((1/i,i)--(1/i,0), gray);
        dot((i,1/i));
        dot((1/i,i));
    }
}
label("$\operatorname{Spec} A$", (1/5,5), N, red);
label("$\operatorname{Spec} B$", (6,0), SW, blue);
label("bad", (0,-5), S);
\end{asy}
\end{adjustbox}
\qquad
\begin{adjustbox}{valign=m}
\begin{asy}
size(0,6cm);
import graph;
real t=0.6;
real f(real x){return 1/x;}
draw(graph(f,0.2,5), red+linewidth(1));
draw(graph(f,-0.2,-5), red+linewidth(1));
draw((-5,-5*t)--(5,5*t), blue+linewidth(1));
for(int i=-4; i<=4; ++i){
    if(i!=0){
        pair P = (i,1/i);
        pair Q = (1/i,i);
        draw(P--foot(P,(0,0),(1,t)), gray);
        draw(Q--foot(Q,(0,0),(1,t)), gray);
        dot(P); dot(Q);
    }
}
label("$\operatorname{Spec} A$", (1/5,5), N, red);
label("$\operatorname{Spec} B$", (5,5*t), NE, blue);
label("good", (0,-5), S);
\end{asy}
\end{adjustbox}
\end{center}
To make the proof work in general, the idea 
is that if the map $\Spec B\to \Spec A$ does not 
look like a covering map, we perturb the variables 
in $A$ a little bit so that it is.
\end{example}
The following proof of this theorem is due to Nagata.
\begin{proof}
    We use induction on $n$. The base case $n=0$
    is obvious. 
    Thus, assume that the statement is true for $n-1$.
    If $I = (0)$, then we are immediately done.

    The strategy is that we will find $z_1,\dots,z_{n-1}\in B$
    such that $y_1,\dots,y_n$ are all integral over ring
    $k[z_1,\dots,z_{n-1}]/J$, which is the subring of $B$ generated by $z_1,\dots,z_{n-1}$ with relations being ideal $J$.
    If we managed to do that, then we would have a chain of ring extension
    $$A=k[x_1,\dots,x_m]\  \stackrel{\text{integral}}\subset\  
    \frac{k[z_1,\dots,z_{n-1}]}J
    \ \stackrel{\text{integral}}\subset\ \frac{k[y_1,\dots,y_n]}I=B,$$
    where the first inclusion is integral extension by 
    applying induction hypothesis on $k[z_1,\dots,z_{m-1}]/J$.
    Thus, by \Cref{prop:integral_transitive},
    we get $B$ is an integral extension of $A$ as desired.

    We now establish such $z_1,\dots,z_{n-1}$.
    Select any nonzero $f\in I$.
    We pick $z_i = y_i - y_n^{r_i}$ where
    $r_1 \gg r_2 \gg \dots \gg r_{m-1} \gg 0$
    (large enough so that $r_i \gg (\deg f) r_{i-1}$).
    $$f(z_1 + y_n^{r_1}, z_2+y_n^{r_2}, 
    \dots, z_{n-1}+y_n^{r_{n-1}}, y_n) = 0.$$

    Imagine expanding this relation.
    All the $y_n$'s will be distinct, and so the largest one
    will be $y_n^N$ for some $N$
    (and other terms, which include some factors of $z_i$,
    have smaller exponent of $y_n$).
    Thus, there exists $\lambda\in k$ such that 
    $\lambda f(z_1 + y_n^{r_1}
    \dots, z_{n-1}+y_n^{r_{n-1}}, y_n)$ is monic in $y_n$.
    Thus, we establish an integral dependence of $y_n$
    in the ring $k[z_1,\dots,z_{n-1}]/\qq$,
    implying that $y_n$ is integral.
    Consequently, $y_i = z_i + y_n^{r_i}$ is integral,
    so we are done.
\end{proof}
\begin{remark}
    If $k$ is infinite, then picking $z_i = y_i-a_iy_n$
    for some $a_1,\dots,a_{n-1}\in k$ also works.
    This is more intuitive because we are doing linear perturbation.
    However, in order to make the proof works for $k$ finite,
    one needs to consider higher-degree perturbation as well.
\end{remark}

\subsection{Nullstellensatz}
We use Noether normalization to prove Hilbert Nullstellensatz,
a theorem that we stated very early in \Cref{chap:rings}.
We first do the weak version.
\begin{theorem}[Weak Nullstellensatz]
    \label{thm:weak_nullstellensatz}
    Let $k$ be a field and $K$ be a field extension
    such that $K$ is finitely-generated as a $k$-algebra.
    Then, $[K:k]$ is a finite extension.
\end{theorem}
\begin{proof}
    By Noether normalization, there exists $m$ such that
    $$\underbrace{k[z_1,\dots,z_m]}_{=A} 
    \ \stackrel{\text{integral}}{\subset}\  K$$
    such that $K$ is an integral extension of $A$.
    However, since $K$ is a field, we get by \Cref{lem:field_case} that $A$ is a field,
    so $m=0$ and $A=k$.
    Thus, $K$ is finitely-generated as $k$-modules.
\end{proof}
\begin{theorem}[Hilbert's Nullstellensatz]
    \label{thm:max_ideal_poly_ring}
    \index{Hilbert's Nullstellensatz}
    Let $k$ be an algebraically-closed field
    and let $A=k[X_1,\dots,X_n]$.
    If $\mathfrak m\subset A$ is maximal,
    then there exists $a_1,\dots,a_n\in k$ such that
    $\mathfrak m = (X_1-a_1,\dots,X_n-a_n)$.
\end{theorem}
\begin{proof}
    We have the map from $k\to A/\mathfrak m$
    by quotient projection.
    By \Cref{thm:weak_nullstellensatz}, 
    $A/\mathfrak m$ is a finite field extension of $k$.
    Since $k$ is algebraically closed,
    we actually get that $A/\mathfrak m = k$.
    Thus, setting $a_i\in k$ be the image of $x_i$ in the
    isomorphism $A/\frakm \to k$,
    we get that $x_i-a_i$ maps to zero,
    so $x_i-a_i\in\mathfrak m$.
    Collating this for all $i$, we get that
    $(x_1-a_1,\dots,x_n-a_n)\subset\mathfrak m$,
    but this ideal is maximal already, so the equality is forced.
\end{proof}
\begin{definition}
Given an ideal $J\subset k[X_1,\dots,X_n]$,
we define the \vocab{vanishing locus}
$$V_m(J) = \{\vec a\in k^n : f(\vec a)=0\text{ for all }f\in J\},$$
which is a closed set (under both Zariski topology and classical topology) of $k^n$.
\notation{$V_m(J)$}{vanishing locus of ideal $J$}
\notation{$I(Z)$}{ideal of functions vanishing at $Z$}

Conversely, given a subset $Z\subseteq k^n$, define the ideal
$$I(Z) = \{f : f(\vec a)=0 \text{ for all }\vec a\in Z\}.$$
\end{definition}
Then, we have the following.
\begin{proposition}
    If $k$ is algebraically closed,
    then the maximal ideals of $k[X_1,\dots,X_n]/I$
    corresponds in bijection to $V_m(I)$.
\end{proposition}
\begin{proof}
    By correspondence theorem,
    maximal ideals of $k[X_1,\dots,X_n]/I$
    correspond to maximal ideals of $k[X_1,\dots,X_n]$
    containing $I$,
    which then correspond to points in $k^n$
    that vanish at every polynomials in $I$.
\end{proof}
In particular, by taking maximal ideals,
we can recover the set of points from ring of function.

Now we investigate the other direction:
given set of points $V_m(I)$, 
can we recover the corresponding ideal $I$?
In other words, we ask of the correspondence
\begin{align*}
    \{\text{ideals of }k[x_1,\dots,x_n]\} &\leftrightarrow \{\text{closed set of }k^n\} \\
    J &\mapsto V_m(J) \\
    I(Z) &\mapsfrom Z.
\end{align*}
is a bijection or not. The answer is not quite,
because ideals $(x^2)$ and $(x)$ have the same
vanishing locus. The fix is that we have to take radical,
and this is strong nullstellensatz.
\begin{theorem}[Strong Hilbert's Nullstellensatz]
    \index{Hilbert Nullstellensatz}
    Let $k$ be an algebraically-closed field
    and $A=k[X_1,\dots,X_n]$. Given an ideal $J\subset A$, we have
	$$I(V_m(J)) = \sqrt J.$$
    To spell out in words, if $f(\vec a)=0$ for all
    $\vec a\in V_m(J)$, then $f^n\in J$ for some $n$.
\end{theorem}
\begin{proof}
    The following proof is known as \textbf{Rabinowitsch's trick}. Let $J$ and $f$ as in the statement.
    
    We are going to add a variable $Y$ and consider an ideal
    $$J' := (J, Yf-1)\ \subseteq\ k[X_1,\dots,X_n,Y].$$
    If $J'\neq k[X_1,\dots,X_n,Y]$, then
    by Krull's theorem \Cref{thm:krull},
    $J'$ is contained in a maximal ideal 
    of $k[X_1\dots,X_n,Y]$, which 
    (by \Cref{thm:max_ideal_poly_ring}) must be of the form
    $(X_1-a_1,\dots,X_n-a_n,b)$ for some $a_1,\dots,a_n,b\in k$.
    The fact that $\mm$ contains $J$ gives
    $(a_1,\dots,a_n)\in V_m(J)$, so
    $f(a_1,\dots,a_n)=0$.
    Moreover, the fact that $\mm$ contais $Yf-1$ 
    gives $f(\vec a)b-1 = 0$, which is a contradiction.
    
    Thus, $J=k[X_1,\dots,X_n,Y]$.
    Therefore, there exists 
    $g_i\in k[X_1,\dots,X_n,Y]$ and $h_i\in J$
    such that
    $$1 =  g_0(fY-1) + \sum_{i=1}^r g_i h_i.$$
    We will now plug in $Y=\frac 1f$,
    working in the localization $k[X_1,\dots,X_n][1/f]$
    to be rigorous. The term outside the summation goes away,
    and we obtain that
    $$1 = \sum_{i=1}^r g_i\left(X_1,\dots,X_n, \frac 1f\right)
    h_i(X_1,\dots,X_n).$$
    Finally, we clear the denominator to get
    $$f^N = \sum_{i=1}^r G_i(X_1,\dots,X_n)
    h_i(X_1,\dots,X_n)\in J,$$
    where $G_i(X_1,\dots,X_n) = f^N g_i(X_1,\dots,X_n,1/f)$
    and $N \gg \deg g_i$ for all $i$,
    so $G_i$ is actually a polynomial.
    Hence, $f\in\sqrt J$, and we are done.
\end{proof}

\begin{remark}
What if $k$ is not algebraically-closed?

If one follows the proof above,
one gets that $A/\mathfrak m$ is a finite extension of $k$,
and it is possible to continue and classify all
maximal ideals of $k[x_1,\dots,x_n]$.
The result is that there is a correspondence between
$$\{\text{maximal ideals of }k[x_1,\dots,x_n]\}
\ \longleftrightarrow\ \frac{\{(a_1,\dots,a_n)\in \overline k^n\}}
{\Gal(\overline k/k)\text{-action}}.$$
In other words, closed points of $k[x_1,\dots,x_n]$
corresponds to Galois orbits of $\overline k^n$.
\Cref{prob:generalized_nullstellensatz} will explore 
this in more detail.


If you take algebraic geometry, you will see that 
a lot of definitions (e.g., morphisms, and when
they are the same) on varieties over $k$
are usually defined by the behavior over $\overline k$.
In the world of affine schemes (i.e., $\Spec A$),
this is completely natural because we naturally
have $\overline k$-points in $\Spec k[x_1,\dots,x_n]$.
\end{remark}

\section{First Steps in Dimension Theory}
\label{sec:dim_intro}
We now use tools developed in this chapter
to study dimension of a variety.
Let us first define dimension.
\begin{definition}[Krull dimension]
    Given any ring $A$, the \vocab{Krull dimension}
    (or just \vocab{dimension})
    of $A$, denoted $\dim A$,
    is the maximum value of $n$ for which
    there exists a chain of prime ideals
    $$\pp_0\subsetneq \pp_1 \subsetneq \dots 
    \subsetneq \pp_n\ \subset \ A$$
    or length $n$.
    \notation{$\dim A$}{dimension of a ring}
\end{definition}
This definition is not too horrible to make because
of the following two reasons.
\begin{itemize}
    \item For any vector space $V$, 
    the dimension of $V$ can be defined as the maximal
    value of $n$ for which there exists a chain of subspaces
    $$V_0\subsetneq V_1\subsetneq \dots\subsetneq V_n 
    \ \subset V$$
    of length $n$. 
    In particular, $V_i$ is an $i$-dimensional subspace.
    This resembles the chain definition given above,
    except that subspaces are replaced by prime ideals.
    \item The chain of prime ideal correspond to 
    a descending chain of (closed) subvarieties
    $$\Spec A/\pp_0 \supsetneq \Spec A/\pp_1 \supsetneq \dots \supsetneq 
    \Spec A/\pp_n.$$
    Intuitively, $\Spec A/\pp_n$ would correspond to a point,
    $\Spec A/\pp_{n-1}$ would correspond to a curve,
    $\Spec A/\pp_{n-2}$ would correspond to a surface, and so on.
    We make $\pp_i$ prime so that one cannot write 
    $\{\text{a line}\} \subsetneq \{\text{union of two lines}\}$
    as an inclusion of a chain, forcing us to increase a dimension.
\end{itemize}
\begin{example}
    \label{ex:dimensions}
    \begin{itemize}
    \item For any field $k$, we have $\dim k=0$.
    \item We claim that $\dim \ZZ=1$ because the prime ideals
    of $\ZZ$ are $(0)$, which is contained in 
    $(2)$, $(3)$, $(5)$, \dots.
    \item More generally, if $A$ is a PID that is not a field, 
    we claim that $\dim A=1$. The ideal $(0)$ is contained 
    in $(a)$ for any non-unit $a\neq 0$.
    Meanwhile, if $(a)$ is prime where $a\neq 0$, 
    then $a$ is irreducible,
    so $(a)$ is maximal, so any nonzero prime ideal is maximal.

    In particular, $\dim k[x]=1$ for any field $k$.
    \item If $k$ is a field, then we can easily show that 
    $\dim k[x_1,\dots,x_n]\geq n$ by exhibiting a chain of 
    prime ideals
    $$(0)\subsetneq (x_1) \subsetneq (x_1,x_2) \subsetneq
    \dots\subsetneq (x_1,\dots,x_n).$$
    On the other hand, showing that $\dim k[x_1,\dots,x_n]\leq n$
    is considerably more difficult because one needs to 
    consider any chain of prime ideals. We will do this in 
    the course of proving \Cref{thm:dim_transcendence}.
    \end{itemize}
\end{example}

Going-up theorem (\Cref{thm:going_up}) controls dimension.
\begin{theorem}[Integral extension has equal dimension]
    \label{thm:dim_integral}
    Let $A$ be a ring and let $B$ be an integral extension
    of $A$ (so $A\subset B$). 
    Then, $\dim A=\dim B$.
\end{theorem}
\begin{proof}
    First, we show $\dim A\geq \dim B$.
    Given any prime ideal chain $\qq_0\subsetneq \qq_1 
    \subsetneq \dots\subsetneq \qq_n$ of $B$,
    one let $\pp_i = \qq_i\cap A$, forming a chain of 
    prime ideal $\pp_0\subsetneq \pp_1 
    \subsetneq \dots\subsetneq \pp_n$ of $A$.
    The equality cannot occur because if $\pp_i=\pp_{i+1}$,
    then $\qq_i\subsetneq \qq_{i+1}$ both lie over $\pp_i$,
    violating incomparability \Cref{thm:lying_over_equal}.

    Next, we show $\dim B\geq \dim A$.
    Given any prime ideal chain $\pp_0\subsetneq \pp_1 
    \subsetneq \dots\subsetneq \pp_n$ of $A$,
    let $\qq_0$ be a prime lying over $\pp_0$
    (exists by \Cref{thm:lying_over}),
    and then by Going Up Theorem (\Cref{thm:going_up}),
    construct the chain $\qq_0\subsetneq \qq_1\subsetneq \dots 
    \subsetneq\qq_n$ such that $\pp_i=\qq_i\cap A$.
    Clearly, $\qq_i\neq \qq_{i+1}$.

    Hence, we conclude that $\dim A=\dim B$.
\end{proof}
Now, we use Noether normalization lemma 
(\Cref{thm:noether_normalization}) to give a neat 
characterization of dimension of a variety.

We first recall that given a field extension
$K/k$, the \vocab{transcendence degree} $\trdeg K/k$ 
is the size of
maximal set $S\subset K$ whose elements are 
$k$-algebraically independent.
For example, 
$$\trdeg k(x_1,\dots,x_n)/k = n$$
because $x_1,\dots,x_n$ are algebraically independent
in $k(x_1,\dots,x_n)$.
\notation{$\trdeg K/k$}{transcendence degree}
\begin{theorem}[Dimension equals transcendence degree]
    \label{thm:dim_transcendence}
    If $k$ is a field, $B$ is a finitely-generated $k$-algebra
    that is a domain
    (i.e., $B=k[x_1,\dots,x_n]/\pp$ for some prime ideal $\pp$),
    and $K=\Frac B$, then 
    $$\dim B = \trdeg K/k.$$
\end{theorem}
\begin{proof}
    We use induction on the transcendence degree $\trdeg K/k=n$.
    Assume that the result is known for transcendence
    degree less than $n$.

    By Noether normalization lemma (\Cref{thm:noether_normalization}),
    there exists $A=k[y_1,\dots,y_m]$ such that 
    $B$ is an integral extension of $A$,
    which means $\dim A=\dim B$ by \Cref{thm:dim_integral}.
    Furthermore, $\Frac B$ is a finite extension of $\Frac A$, 
    so $m = \trdeg(\Frac(A)/k) =\trdeg(\Frac(B)/k)= n$.

    Thus, it suffices to show that $\dim k[y_1,\dots,y_n]=n$.
    We have shown in \Cref{ex:dimensions} above that 
    the dimension is at least $n$.
    To show that it is at most $n$, suppose for contradiction
    that there is a chain of length $n+1$:
    $$\pp_0 \subsetneq \pp_1 \subsetneq \dots \subsetneq \pp_{n+1}.$$
    Replace $\pp_0$ with $(0)$. Choose an element $g\in \pp_1$,
    let $f$ be an irreducible factor of $g$ lying in $\pp_1$,
    and replace $\pp_1$ with $(f)$. (Here, we use the fact that
    $k[x_1,\dots,x_n]$ is a UFD to see  that $(f)$ is prime.)
    Thus, the above chain simplifies to
    $$(0) \subsetneq (f)\subsetneq \pp_2
    \subsetneq \dots \subsetneq \pp_{n+1},$$
    which induces the chain of prime ideal of length $n$ on ring 
    $C = k[x_1,\dots,x_n]/(f)$.

    However, note that $\Frac C = k(x_1,\dots,x_n)/(f)$,
    which has transcendence degree $n-1$.
    Therefore, by induction hypothesis, $\dim C = n-1$,
    contradicting the above paragraph.
\end{proof}
\begin{example}
    If $f$ is any irreducible polynomial in $k[x_1,\dots,x_n]$,
    then $\dim k[x_1,\dots,x_n]/(f) = n-1$.
\end{example}
This covers dimension theory in the case of 
finitely-generated algebra over $k$. 
However, we still don't know how to prove, for example,
that $\dim \ZZ[x_1,\dots,x_n] = n+1$.
We will develop a more robust dimension theory in 
\Cref{chap:dimension} that will allow us to prove that.
%
%
\section{Problems}
\begin{problem}
Decide whether each of the following ring is an integral extension of $\mathbb C[x]$.
\begin{alphenum}
\item $\mathbb C[x,y]/(x^2-y^2)$.
\item $\mathbb C[x,y]/(xy)$.
\item $\mathbb C[x,y]/(xy-1)$.
\end{alphenum}
Do your geometric intuition and algebraic proof match?
\end{problem}
\begin{problem}
If $B$ is an integral extension of $A$ and $\qq\subset B$
is a prime ideal lying above $\pp\subset A$,
is it necessarily true that $B_\qq$ is integral over $A_\pp$?
Compare with \Cref{prop:integral_quotient_localize} (b).
\end{problem}
\begin{problem}
Let $d$ be a squarefree integer. 
Prove that the integral closure of $\ZZ$ in $\QQ(\sqrt d)$ is 
$$\begin{cases}
    \ZZ[\sqrt d] & \text{if } d\equiv 2,3\pmod 4 \\
    \ZZ\left[\tfrac{1+\sqrt d}2\right] & 
    \text{if } d\equiv 1\pmod 4.
\end{cases}$$
\end{problem}
\begin{problem}
Prove that the intergral closure of $\mathbb C[x,y]/(y^2-x^3-x^2)$
is isomorphic to $\mathbb C[t]$, where $t=\frac yx$.
Can you draw picture similar to \Cref{ex:integral_closure}?
\end{problem}
\begin{problem} 
\label{prob:integrally_closed}
Prove that $\mathbb C[x,y,z]/(x^2-yz)$ is integrally closed.%
\footnote{Note that the surface $x^2-yz=0$ has a singularity at $(0,0,0)$.}

(Hint: minimal polynomial test, \Cref{lem:min_poly_test}.)
\end{problem}
\begin{problem} 
Let $K$ be a finite, Galois extension of $\mathbb Q$. Let $\mathcal O_K$ be the ring of integer of $K$. Let $p$ be a rational prime, and $\mathfrak p$ and $\mathfrak q$ are two prime ideals of $\mathcal O_K$ lying above $(p)$. Prove that there exists $\sigma\in\operatorname{Gal}(K/\mathbb Q)$ such that $\sigma\mathfrak p = \mathfrak q$.

(Hint: use prime avoidance, \Cref{prob:prime_avoidance}.)
\end{problem}
\begin{problem}
\label{prob:going_down}
Using the images shown below,
come up with counterexamples to show that 
the Going-down theorem is not true if 
we remove the hypothesis that 
\begin{alphenum}
\item $B$ is a domain.
\begin{center}
\begin{asy}
size(4.5cm);
dot("$\mathfrak q'$", (0,0.75), N, blue);
draw((-1,0.5)--(1,0.5), red);
label("\Large$\Downarrow$",(0,0));
draw((-1,-0.5)--(1,-0.5), red);
dot("$\mathfrak q$", (0,-0.5), N, blue);
label("$\mathfrak p$", (-1,-0.5), NW, red);
\end{asy}
\end{center}
\item $A$ is integrally closed.
\end{alphenum}
\begin{center}
\begin{adjustbox}{valign=m}
\begin{asy}
size(0,5cm);
fill(box((-1,-1.5), (1,1.5)), 0.2*black+0.8*white);
draw((-1,0.6)--(1,0.6), dashed);
draw((-1,-0.6)--(1,-0.6), dashed);
label("glue", (-1.5,0));
draw((-1.5,0.2)--(-1,0.6), EndArrow);
draw((-1.5,-0.2)--(-1,-0.6), EndArrow);
draw((-1,-1)--(1,1), red);
dot((-0.6,-0.6), blue+4);
dot((0.6,0.6), blue+4);
dot("$\mathfrak q'$", (-0.6,0.6), NE, deepgreen+4);
\end{asy}
\end{adjustbox}
\qquad
\begin{adjustbox}{valign=m}
{\Large $\Rightarrow$}
\end{adjustbox}
\qquad
\begin{adjustbox}{valign=m}
\begin{asy}
settings.tex="pdflatex";
settings.prc=false;
settings.render=0;
import three;
import graph3;
size(0,5cm);

currentprojection = perspective(-1,-4,1.2);
real f(pair z) {return sqrt(z.x^3 + z.x^2);}
real g(pair z) {return -sqrt(z.x^3 + z.x^2);}


real x(real t) {return t^2-1;}
real y(real t) {return 0.5*t;}
real z(real t) {return t*(t^2-1);}

surface S1 = surface(f,(-1,-1), (0,1), nx=20, ny=20);
surface S2 = surface(g,(-1,-1), (0,1), nx=20, ny=20);
surface S3 = surface(f,(0,-1), (0.67,1), nx=20, ny=20);
surface S4 = surface(g,(0,-1), (0.67,1), nx=20, ny=20);

draw(S3, gray);
draw(S4, gray);
draw(S2, gray);
draw(graph(x,y,z,0,1.3,operator ..), red);
draw(S1, gray);
draw(graph(x,y,z,-1.3,0,operator ..), red);


dot((0,0.5,0), blue);
dot("$\mathfrak q$",(0,-0.5,0), blue);
label("$\mathfrak p$",(x(1.32), y(1.32), z(1.32)), red);
\end{asy}
\end{adjustbox}
\end{center}
\end{problem}
\begin{problem}
Find a Noether normalization and the dimension of the ring
$$A = \frac{\mathbb C[x,y,z]}{(x^2-y^3-1, xyz-1)}.$$
\end{problem}
\begin{problem}
\label{prob:generalized_nullstellensatz}
Let $\mathfrak n$ be a maximal ideal of $A=\mathbb R[x_1,\dots,x_n]$.
\begin{alphenum}
\item Prove that there exists $a=(a_1,\dots,a_n)\in\mathbb C^n$ such that if $\mathfrak m_a = (x_1-a_1,\dots,x_n-a_n)\subset\mathbb C[x_1,\dots,x_n]$, then $\mathfrak m_a\cap A = \mathfrak n$.
\item Let $a,b\in\mathbb C^n$. Prove that $\mathfrak m_a \cap A= \mathfrak m_b\cap A$ if and only if $a=b$ or $a=\overline b$ (where $\overline b$ is the complex conjugate of $b$.)
\end{alphenum}
Thus, we establish a bijection between maximal ideals of $\mathbb R[x_1,\dots,x_n]$ and tuples in $\mathbb C^n$ modulo conjugation. 
You can generalize this result by replacing $\RR$ and $\CC$
to $k$ and $\overline k$, respectively.
\end{problem}
\begin{problem}
Let $k$ be a field and $A$ and $B$ are finitely-generated $k$-algebra.
\begin{alphenum}
\item Assume that $k$ is algebraically closed.
Prove that there exists a bijection between 
$$\{\text{maximal ideals of }A\otimes_k B\} 
\longleftrightarrow 
\{\text{maximal ideals of }A\}
\times \{\text{maximal ideals of }B\}.$$
Geometric interpretation: $A\otimes_k B$
is the ring corresponding to product of varieties.
\item Prove that $\dim(A\otimes_k B) = \dim A + \dim B$,
where $\dim$ in both sides are dimension as rings.
\end{alphenum}
\end{problem}
\begin{problem}
Let $k$ be a field, and let $A$ be a finitely-generated $k$-algebra that is a domain.
\begin{alphenum}
\item Let $B$ be an integral extension of $A$ that is a domain. Suppose that $\mathfrak p_1 \subsetneq  \mathfrak p_2 \subsetneq \mathfrak p_3 \subset A$ is a chain of prime ideals in $A$. Let $\mathfrak q_1\subsetneq\mathfrak q_3\subset B$ be prime ideals above $\mathfrak p_1$ and $\mathfrak p_3$, respectively. Prove that there exists a prime ideal $\mathfrak q_2$ such that $\mathfrak q_1\subsetneq \mathfrak q_2\subsetneq \mathfrak q_3$.
Note that $\qq_2$ need not lie over $\pp_2$.

Hint: Without loss of generality, 
assume $\mathfrak p_1 = (0)$ and $\mathfrak q_1=(0)$. 
Why can we do this?
\item Prove that all maximal chains of prime ideals in $A$ have length exactly $\dim A$. 
(Here, maximal is in the sense that one cannot insert another prime ideal at the boundary or 
between any two consecutive ideals in the chain.)
\item Prove that for all prime ideals $\mathfrak p \subsetneq \mathfrak q$ of $A$, the length of all maximal chains of prime ideals starting at $\mathfrak p$ and ending at $\mathfrak q$ are the same. (Rings with this property are called \textbf{catenary}.)
\end{alphenum}
\end{problem}